\subsection{A coefficient lemma for signed traces}

The preceding decomposition leads to the following abstract statement.
The operator $E$ describes the exterior, $C$ contains the central
interval, and the two word families collect the perturbations at its
ends. Their lengths include the unperturbed blocks between chosen
perturbations. The lemma extracts the coefficient of each peripheral
eigenvalue from a difference of traces.

\begin{lemma}[Signed trace coefficients]\label{j:trace-extraction}
Let $E$ be a Hilbert--Schmidt operator on a complex Hilbert space with
spectral radius one. Let $\mathcal P=\{\zeta\in\sigma(E):|\zeta|=1\}$.
For each $\zeta\in\mathcal P$ let $P_\zeta$ be its Riesz projection,
$d_\zeta=\operatorname{rank}P_\zeta$, and
$J_\zeta=(E-\zeta I)P_\zeta$, with
$J_\zeta^{q_\zeta}=0$ on $\operatorname{ran}P_\zeta$.
Let $C$ be trace class. Suppose two countable families of bounded
operators $W_+,W_-$ have integer lengths $\ell(W_\pm)\ge0$ and, for
some $\epsilon>0$,
\begin{equation}\label{j:word-moment}
 \sum_{W_\pm}e^{\epsilon\ell(W_\pm)}\|W_\pm\|<\infty.
\end{equation}
For integers $s\ge0$ put
\[
 F_s=\sum_{\ell(W_+)+\ell(W_-)\le s}
 \operatorname{tr}(CW_+E^{s-\ell(W_+)-\ell(W_-)}W_-).
\]
Assume that $a\ge0$ is an integer and $R\in\mathbb C$ satisfy
$F_s-R\operatorname{tr}E^{s+a}\to0$ as $s\to\infty$.
Only sufficiently large $s$, with $s+a\ge2$, are used in this limit.
Then the coefficientwise absolutely convergent polynomial
\[
 \mathcal P_\zeta(z)=\sum_{W_+,W_-}\zeta^{-\ell}
 \operatorname{tr}\left(CW_+
 \left[\sum_{t=0}^{q_\zeta-1}
 \binom{z-\ell}{t}\zeta^{-t}J_\zeta^tP_\zeta\right]W_-\right),
 \quad \ell=\ell(W_+)+\ell(W_-),
\]
is identically $R d_\zeta\zeta^a$. In particular,
\[
 R=\frac{\mathcal P_\zeta(0)}{d_\zeta\zeta^a},\qquad
 |R|\le K\|C\|_1,
\]
where $K$ depends on $E$ and the bounds in \eqref{j:word-moment}, but
not on $C$ or $a$. More precisely it may be chosen in terms of
$\epsilon$, the two weighted word sums, the finite nilpotent orders
$q_\zeta$, and the norms $\|J_\zeta^tP_\zeta\|$.
Thus the conclusion is uniform for a family with uniform bounds on
these data. No peripheral semisimplicity is assumed.
\end{lemma}

\begin{proof}
Compactness makes $\mathcal P$ finite and nonempty. On its finite Riesz
blocks,
\[
 E^nP_\zeta=\zeta^n\sum_{t=0}^{q_\zeta-1}
 \binom nt\zeta^{-t}J_\zeta^tP_\zeta.
\]
The complementary restriction $E_<$ has spectral radius strictly less
than one. Thus the remainder in this expansion is $O(\rho^n)$ in
operator norm for some $\rho<1$, and
$\|E^n\|\le K_0(1+n)^q$ for a fixed $q$.
For $n\ge2$, the complementary trace norm is bounded by
$\|E_<\|_2^2\|E_<^{n-2}\|$. Every positive nilpotent power has trace zero
on a finite Riesz block, so
\begin{equation}\label{j:bare-trace-expansion}
 \operatorname{tr}E^n=\sum_{\zeta\in\mathcal P}
 d_\zeta\zeta^n+O(\rho_1^n),\qquad \rho_1<1.
\end{equation}

The inequality $|\operatorname{tr}(CAB)|\le\|C\|_1\|A\|\|B\|$
and the polynomial power bound prove absolute convergence of $F_s$.
A coefficient of $\binom{z-\ell}{t}$ is bounded by a constant times
$(1+\ell)^t$, so the exponential moments in \eqref{j:word-moment}
also prove coefficientwise convergence of $\mathcal P_\zeta$.
For pairs with $\ell\le s/2$, replacing the middle power by its
peripheral expansion costs $O(\rho^{s/2})\|C\|_1$. For the remaining
pairs use the polynomial bound and the exponential tail of the word
sums. Their total is at most a polynomial in $s$ times
$e^{-\epsilon s/2}\|C\|_1$. The same estimate permits extending the
peripheral sum to all word pairs, using the polynomial growth of the
generalized binomial coefficients. Hence
\[
 F_s=\sum_{\zeta\in\mathcal P}\zeta^s\mathcal P_\zeta(s)+o(1).
\]
Together with the assumed difference of traces and
\eqref{j:bare-trace-expansion}, this gives
\[
 \sum_{\zeta\in\mathcal P}\zeta^s
 \big(\mathcal P_\zeta(s)-Rd_\zeta\zeta^a\big)\longrightarrow0.
\]
Distinct unit-modulus exponentials cannot cancel a nonzero polynomial
coefficient in this way. Divide by the largest occurring power $s^d$,
multiply by $\overline\eta^{s}$ for one $\eta\in\mathcal P$, and take
Ces\`aro means. The geometric-series formula makes every frequency
except $\eta$ vanish, leaving its leading coefficient. That coefficient
is zero; descend through the degrees. This proves the polynomial
identity. Evaluation at zero is bounded by
\[
 K_1\|C\|_1\sum_{W_+,W_-}
 \|W_+\|\|W_-\|(1+\ell)^{q_\zeta-1},
\]
which is finite by \eqref{j:word-moment}. Since $d_\zeta\ge1$ and
$|\zeta^a|=1$, the asserted bound follows.
\end{proof}
